% ***************************************************************************************************
%
%	Szablon pracy magisterskiej dla Politechniki Wrocławskiej w wersji dwustronnej.
%	Autor:	Tomasz Strzałka
%
% ***************************************************************************************************

% Styl dwustronny z domyślną wielkością czcionki 10pt oraz oddzieloną stroną tytułową (titlepage).
% Domyślnie rodziały rozpoczynają się na stronie prawej (openright).
\documentclass{book}

% ***************************************************************************************************
% Ustawienia języka
% ***************************************************************************************************

% Podstawowe ustawienia języka, według którego formatowany będzie dokument
\usepackage[english]{babel}
\usepackage{csquotes}
% Pakiet babel dla polskiego języka powoduje konflikt z pakietem amssymb.
% Polecenie '\lll' definiują oba pakiety - porządana jest druga definicja.

% W przypadku wielojęzykowości ustawia główny język dokumentu
\selectlanguage{english}

% Kodowanie dokumentu
\usepackage[utf8]{inputenc}

% Dowolny rozmiar czcionek, kodowanie znaków
\usepackage{lmodern}

% Polskie wcięcia akapitów
% \usepackage{indentfirst}

% Polskie łamanie wyrazów
% \usepackage[plmath]{english}

% Przecinek w wyrażeniach matematycznych zamiast kropki
% \usepackage{icomma}

% Polskie formatowanie typograficzne
% \frenchspacing

% Zapewnia liczne usprawnienia wyświetlania i organizacji matematycznych formuł. 
\usepackage{amsmath}

% Wprowadza rozszerzony zestaw symboli m.in. \leadsto
\usepackage{amssymb}

% Dodatkowa, ,,kręcona'' czcionka matematyczna
\usepackage{mathrsfs}

% Dodatkowe wsparcie dla środowiska mathbb, które nie wspiera domyślnie cyfr (\mathbb{})
\usepackage{bbold}

% Fixes/improves amsmath
\usepackage{mathtools}



% ***************************************************************************************************
% Kolory  
% ***************************************************************************************************

% Umożliwia kolorowanie poszczególnych komórek tabeli
\usepackage[table]{xcolor}% http://ctan.org/pkg/

% Umożliwia łatwą zmianę koloru linii w tabeli
\usepackage{tabu}

% Umożliwia rozszerzoną kontrolę nad kolorami.
\usepackage{xcolor}

% Definicje kolorów
\definecolor{lgray}{HTML}{9F9F9F}
\definecolor{dgray}{HTML}{5F5F5F}
% lgray				-	nazwa nowo zdefiniowanego koloru
% HTML				-	model kolorów
% CCCCCC			-	wartość koloru zgodna z modelem

% ***************************************************************************************************
% Algorytmy 
% ***************************************************************************************************

% Udostępnia środowisko do konstruowania pseudokodów
% \usepackage[ruled,vlined,linesnumbered,longend,algochapter]{algorithm2e}
% ruled	- poziome kreski na początku i końcu algorytmu, podpis na górze oddzielony również kreską poziomą
% vlined - pionowe kreski łączące początek polecenia z jego końcem
% linesnumbered	- numerowanie kolejnych wierszy algorytmu
% longend - długie końcówki np. ifend, forend itd.
% algochapter - numeracja z rozdziałami

% Zmiana rozmiaru komentarzy
% \newcommand\algcomment[1]{
	% \footnotesize{#1}
% }

% Ustawienie zadanego stylu dla komentarzy
% \SetCommentSty{algcomment}

% Wyśrodkowana tylda
\usepackage{textcomp}%
\newcommand{\textapprox}{\raisebox{0.5ex}{\texttildelow}}

% Listowanie kodów źródłowych
\usepackage{listings} 

% ***************************************************************************************************
% Marginesy 
% ***************************************************************************************************

% Ustawienia rozmiarów stron i ich marginesów
\usepackage[headheight=18pt, top=25mm, bottom=25mm, left=25mm, right=25mm]{geometry}
% headheight		-	wysokość tytułów
% top				-	margines górny
% bottom			-	margines dolny
% left				-	margines lewy
% right				-	margines prawy

% Usunięcie górnego marginesu dla środowisk
\makeatletter
\setlength\@fptop{0\p@}	
\makeatother

% ***************************************************************************************************
% Styl 
% ***************************************************************************************************

% Definiuje środowisko 'titlingpage', które zapewnia pełną kontrolę nad układem strony tytułowej.
\usepackage{titling}


% Umożliwia modyfikowanie stylu spisu treści
\usepackage{tocloft}	

\tocloftpagestyle{tableOfContentStyle}

% Definiowanie własnych stylów nagłówków i/lub stopek
\usepackage{fancyhdr}

% Domyślny styl dla pracy 
\fancypagestyle{custom}{
	\fancyhf{}									% wyczyść stopki i nagłówki
	\fancyhead[RO]{								% Prawy, nieparzysty nagłówek
		\hrulefill \hspace{16pt} \large Chapter \thechapter
		\put(-472.1, 12.1){%
			\makebox(0,0)[l]{%
				\includegraphics[width=0.05\textwidth]{Figures/pwr-logo}
			}
		}
		\put(-443,5.5){%
			\makebox(0,0)[l]{%
				\small Wrocław University of Science and Technology
			}
		}
	}
	\fancyhead[LE]{								% Lewy, parzysty nagłówek
		\large Chapter \thechapter \hspace{16pt} \hrulefill 
		\put(-22, 12.1){%
			\makebox(0,0)[l]{%
				\includegraphics[width=0.05\textwidth]{Figures/wppt-logo}
			}
		}
		\put(-220,5.5){%
			\makebox(0,0)[l]{%
				\small Faculty of Fundamental Problems of Technology
			}
		}
	}
	\fancyfoot[LE,RO]{							% Stopki
		\thepage
	}
	\renewcommand{\headrulewidth}{0pt}			% Grubość linii w nagłówku
	\renewcommand{\footrulewidth}{0.2pt}		% Grubość linii w stopce
}


% Domyślny styl dla bibliografii
\fancypagestyle{bibliographyStyle}{
	\fancyhf{}									% wyczyść stopki i nagłówki
	\fancyhead[RO]{								% Prawy, nieparzysty nagłówek
		\hrulefill \hspace{16pt} \large Appendix \thechapter
		\put(-472.1, 12.1){%
			\makebox(0,0)[l]{%
				\includegraphics[width=0.05\textwidth]{Figures/pwr-logo}
			}
		}
		\put(-443,5.5){%
			\makebox(0,0)[l]{%
				\small Wrocław University of Science and Technology
			}
		}
	}
	\fancyhead[LE]{								% Lewy, parzysty nagłówek
		\large Bibliography \hspace{16pt} \hrulefill 
		\put(-22, 12.1){%
			\makebox(0,0)[l]{%
				\includegraphics[width=0.05\textwidth]{Figures/wppt-logo}
			}
		}
		\put(-220,5.5){%
			\makebox(0,0)[l]{%
				\small Faculty of Fundamental Problems of Technology
			}
		}
	}
	\fancyfoot[LE,RO]{							% Stopki
		\thepage
	}
	\renewcommand{\headrulewidth}{0pt}			% Grubość linii w nagłówku
	\renewcommand{\footrulewidth}{0.2pt}		% Grubość linii w stopce
}

% Domyślny styl dla dodatków
\fancypagestyle{appendixStyle}{
	\fancyhf{}									% wyczyść stopki i nagłówki
	\fancyhead[RO]{								% Prawy, nieparzysty nagłówek
		\hrulefill \hspace{16pt} \large Appendix \thechapter
		\put(-472.1, 12.1){%
			\makebox(0,0)[l]{%
				\includegraphics[width=0.05\textwidth]{Figures/pwr-logo}
			}
		}
		\put(-443,5.5){%
			\makebox(0,0)[l]{%
				\small Wrocław University of Science and Technology
			}
		}
	}
	\fancyhead[LE]{								% Lewy, parzysty nagłówek
		\large Dodatek \thechapter \hspace{16pt} \hrulefill 
		\put(-22, 12.1){%
			\makebox(0,0)[l]{%
				\includegraphics[width=0.05\textwidth]{Figures/wppt-logo}
			}
		}
		\put(-210,5.5){%
			\makebox(0,0)[l]{%
				\small Faculty of Fundamental Problems of Technology
			}
		}
	}
	\fancyfoot[LE,RO]{							% Stopki
		\thepage
	}
	\renewcommand{\headrulewidth}{0pt}			% Grubość linii w nagłówku
	\renewcommand{\footrulewidth}{0.2pt}		% Grubość linii w stopce
}

% Osobny styl dla stron zaczynających rozdział/spis treści itd. (domyślnie formatowane jako "plain")
\fancypagestyle{chapterBeginStyle}{
	\fancyhf{}%
	\fancyfoot[LE,RO]{
		\thepage
	}
	\renewcommand{\headrulewidth}{0pt}
	\renewcommand{\footrulewidth}{0.2pt}
}

% Styl dla pozostałych stron spisu treści
\fancypagestyle{tableOfContentStyle}{
	\fancyhf{}%
	\fancyfoot[LE,RO]{
		\thepage
	}
	\renewcommand{\headrulewidth}{0pt}
	\renewcommand{\footrulewidth}{0.2pt}
}

% Formatowanie tytułów rozdziałów i/lub sekcji
\usepackage{titlesec}

% Formatowanie tytułów rozdziałów
\titleformat{\chapter}[hang]					% kształt
{
	\vspace{-10ex}
	\Huge
	\bfseries
}												% formatowanie tekstu modyfikowanego elementu
{}												% etykieta występująca przed tekstem modyfikowanego elementu, niewidoczna w spisie treści
{
	10pt
}												% odstęp formatowanego tytułu od lewego marginesu/etykiety
{
	\Huge
	\bfseries
}												% formatowanie elementów przed modyfikowanym tytułem
[
\vspace{2ex}
%\rule{\textwidth}{0.4pt}
%\vspace{-4ex}
]												% dodatkowe formatowanie stosowane poniżej modyfikowanego tytułu


% Formatowanie tytułów sekcji
\titleformat{\section}[hang]					% kształt
{
	\vspace{2ex}
%	\titlerule\vspace{1ex}
	\Large\bfseries
}												% formatowanie tekstu modyfikowanego elementu
{
	\thesection									% etykieta występująca przed tekstem modyfikowanego elementu, niewidoczna w spisie treści
}
{
	0pt
}												% odstęp formatowanego tytułu od lewego marginesu/etykiety
{
	\Large
	\bfseries
}												% formatowanie elementów przed modyfikowanym tytułem

% ***************************************************************************************************
% Linki
% ***************************************************************************************************

% Umożliwia wstawianie hiperłączy do dokumentu
\usepackage{hyperref}							% Aktywuje linki

\hypersetup{
	colorlinks	=	true,					% Koloruje tekst zamiast tworzyć ramki.
	linkcolor		=	blue,					% Kolory: referencji,
        citecolor		=	blue,					% cytowań,
	urlcolor		=	blue					% hiperlinków.
}

% Do stworzenia hiperłączy zostanie użyta ta sama (same) czcionka co dla reszty dokumentu
\urlstyle{same}




% ***************************************************************************************************
% Linki
% ***************************************************************************************************

% Umożliwia zdefiniowanie własnego stylu wyliczeniowego
\usepackage{enumitem}

% Nowa lista numerowana z trzema poziomami
\newlist{myitemize}{itemize}{3}

% Definicja wyglądu znacznika pierwszego poziomu
\setlist[myitemize,1]{
	label		=	\textbullet,
	leftmargin	=	4mm}

% Definicja wyglądu znacznika drugiego poziomu
\setlist[myitemize,2]{
	label		=	$\diamond$,
	leftmargin	=	8mm}

% Definicja wyglądu znacznika trzeciego poziomu
\setlist[myitemize,3]{
	label		=	$\diamond$,
	leftmargin	=	12mm
}


% ****************************** Biblatex ***********************************************************
\usepackage[
    backend=biber,
    style=numeric,
	sorting=none
  ]{biblatex}
\addbibresource{qliom.bib}

% ***************************************************************************************************


% ***************************************************************************************************
% Inne pakiety
% ***************************************************************************************************

% Dołączanie rysunków
\usepackage{graphicx}

% Figury i przypisy
\usepackage{caption}
\usepackage{subcaption}

% Umożliwia tworzenie przypisów wewnątrz środowisk
\usepackage{footnote}

% Umożliwia tworzenie struktur katalogów
\usepackage{dirtree}

% Rozciąganie komórek tabeli na wiele wierszy
\usepackage{multirow}

% Precyzyjne obliczenia szerokości/wysokości dowolnego fragmentu wygenerowanego przez LaTeX
\usepackage{calc}

% my packages
\usepackage{physics}
\usepackage{microtype}
\usepackage{cancel}
\usepackage{color}
\renewcommand{\CancelColor}{\color{red}}
% ***************************************************************************************************
% Matematyczne skróty
% ***************************************************************************************************

% Skrócony symbol liczb rzeczywistych
\newcommand{\RR}{\mathbb{R}}

% Skrócony symbol liczb naturalnych
\newcommand{\NN}{\mathbb{N}}

% Skrócony symbol liczb wymiernych
\newcommand{\QQ}{\mathbb{Q}}

% Skrócony symbol liczb całkowitych
\newcommand{\ZZ}{\mathbb{Z}}

% Skrócony symbol logicznej implikacji
\newcommand{\IMP}{\rightarrow}

% Skrócony symbol  logicznej równoważności
\newcommand{\IFF}{\leftrightarrow}

% Symbole operatorów spinowych
\newcommand{\Id}{\mathcal{I}}
\newcommand{\Sp}{S^{+}}
\newcommand{\Sm}{S^{-}}
\newcommand{\Sz}{S^{z}}
\newcommand{\Sx}{S^{x}}
\newcommand{\Sy}{S^{y}}

% ***************************************************************************************************
% Środowiska
% ***************************************************************************************************

% Środowisko do twierdzeń
\newtheorem{theorem}{Theorem}[chapter]

% Środowisko do lematów
\newtheorem{lemma}{Lemma}[chapter]

% Środowisko do przykładów
\newtheorem{example}{Example}[chapter]

% Środowisko do wniosków
\newtheorem{corollary}{Corollary}[chapter]

% Środowisko do definicji
\newtheorem{definition}{Definition}[chapter]

% Środowisko do dowodów
\newenvironment{proof}{
	\par\noindent \textbf{Proof.}
}{
\begin{flushright}
	\vspace*{-6mm}\mbox{$\blacklozenge$}
\end{flushright}
}

% Środowisko do uwag
\newenvironment{remark}{
	\bigskip \par\noindent \small \textbf{Remark.}
}{
\begin{small}
	\vspace*{4mm}
\end{small}
}


% ***************************************************************************************************
% Dokument
% ***************************************************************************************************

\frontmatter

\begin{document}

	\begin{titlingpage}
		\vspace*{\fill}
		\begin{center}
			\begin{picture}(300,510)
				\put(11,520){\makebox(0,0)[l]{\large \textsc{Faculty of Fundamental Problems of Technology}}}
				\put(11,500){\makebox(0,0)[l]{\large \textsc{Wrocław University of Science and Technology}}}
% Tytuł pracy
				\put(80,320){\Huge \textsc{Non-commuting}}
				\put(80,280){\Huge \textsc{integrals of motion}}
				\put(80,240){\Huge \textsc{in XXZ model}}
% Autor pracy
				\put(90,200){\makebox(0,0)[l]{\large \textsc{Jakub Pawłowski}}}
				\put(90,180){\makebox(0,0)[l]{\large \textsc{Index number: 250193}}}

				\put(200,100){\makebox(0,0)[l]{\large Bachelor thesis}}
				\put(200,80){\makebox(0,0)[l]{\large under supervision of}}
% dane promotora
				\put(200,60){\makebox(0,0)[l]{\large prof. dr hab. Marcin Mierzejewski}}
				
				\put(115,-70){\includegraphics[width=0.15\textwidth]{Figures/pwr}}
				\put(106,-80){\makebox(0,0)[bl]{\large \textsc{Wrocław 2021}}}
			\end{picture}
		\end{center}	
		\vspace*{\fill}
	\end{titlingpage}
	
        \cleardoublepage{}
		
	\pagenumbering{Roman}
	\pagestyle{tableOfContentStyle}
	\tableofcontents
	\cleardoublepage{}
		
	% ***************************************************************************************************
	% Wstęp
	% ***************************************************************************************************
	
	\pagestyle{custom}
	\mainmatter{}
	
	% ***************************************************************************************************
	% Rodziały
	% ***************************************************************************************************

	\chapter{Introduction}
\thispagestyle{chapterBeginStyle}

In this work we will describe \ldots~\cite{Mierzejewski2015b}.

	\cleardoublepage{}

	\chapter{Energy current}
\thispagestyle{chapterBeginStyle}
\label{chap:energy current}

\paragraph{}In order to test our QLIOM finding algorithm and the correctness of its implementation, we investigate the known case of
energy current in Spin\(-1/2\) XXZ model~\autocite*{Mierzejewski2015Approx}. For the sake of completeness, derivation of
spin energy current for the general XYZ model will be presented, following the definitions in \textcite{Zotos1997}.
We start with the general XYZ Hamiltonian with periodic boundary conditions:
\begin{equation}
    H_{XYZ} = \sum_{i=1}^L  \left( J_{x} \Sx_{i}\Sx_{i+1} + J_{x} \Sy_{i}\Sy_{i+1} + J_{z} \Sz_{i}\Sz_{i+1} \right)
\end{equation}
It is easy to see that this Hamiltonian can be represented as a sum of operators supported on two consecutive sites:
\begin{equation}
    H_{XYZ} = \sum_{i=1}^L h_{i,i+1}
\end{equation}
where \(h_{i,i+1} = J_{x} \Sx_{i}\Sx_{i+1} + J_{x} \Sy_{i}\Sy_{i+1} + J_{z} \Sz_{i}\Sz_{i+1} \) and periodic boundary conditions
require that \(h_{L,L+1} = h_{L,1}\). The energy operator is a conserved quantity, thus the time evolution of its local density
is given by the discrete continuity equation:
\begin{equation}
    \dv{h_{i,i+1}(t)}{t} + \div{j_{i}^{E}(t)} = 0 
    \label{eq:discrete continuity}
\end{equation}
where \(\div{j_{i}^E(t)} \equiv j_{i+1}^E(t) - j_{i}^E(t)\) is the discrete divergence of local energy current and \(h_{i,i+1}(t) = e^{i H_{XYZ}t} h_{i,i+1} e^{-i H_{XYZ} t}\).
On the other hand, time evolution of an arbitrary operator is determined
by the Heisenberg equations:
\begin{equation}
    \dv{h_{i,i+1}(t)}{t} = i \comm{H_{XYZ}}{h_{i,i+1}(t)}
    \label{eq:heisenberg equation}
\end{equation}
Combining equations~\eqref{eq:discrete continuity} and~\eqref{eq:heisenberg equation} we obtain the defining equations for
the spin energy current:
\begin{equation}
    j_{i+1}^E - j_{i}^E = - i \comm{H_{XYZ}}{h_{i,i+1}} = i \comm{h_{i,i+1}}{H_{XYZ}} = i \sum_{k=1}^L \comm{h_{i,i+1}}{h_{k,k+1}} 
    \label{eq:energy current defining equation}
\end{equation}
Similar equations can be written for any operator being a sum of local operators (\textcolor{blue}{with support 2?}) such as
the total spin operator or particle number operator in Hubbard model. Equation~\eqref{eq:energy current defining equation} is
conceptually simple, yet quite tedious to solve due to the amount of commutators present. Luckily, leveraging commutator properties
to our advantage will allow us to simplify the calculations. Let us begin with inserting the definition of \(h_{i,i+1}\) into 
equation~\eqref{eq:energy current defining equation}:
\begin{align*}
    \comm{h_{i,i+1}}{h_{k,k+1}} =& \comm{J_{x} \Sx_{i}\Sx_{i+1} + J_{x} \Sy_{i}\Sy_{i+1} + J_{z} \Sz_{i}\Sz_{i+1}}{J_{x} \Sx_{k}\Sx_{k+1} + J_{x} \Sy_{k}\Sy_{k+1} + J_{z} \Sz_{k}\Sz_{k+1}}\\
    =& J_{x} J_{y}\comm{\Sx_{i}\Sx_{i+1}}{\Sy_{k}\Sy_{k+1}} + J_{x} J_{z}\comm{\Sx_{i}\Sx_{i+1}}{\Sz_{k}\Sz_{k+1}} + J_{y} J_{x}\comm{\Sy_{i}\Sy_{i+1}}{\Sx_{k}\Sx_{k+1}}\\
    +& J_{y} J_{z}\comm{\Sy_{i}\Sy_{i+1}}{\Sz_{k}\Sz_{k+1}} + J_{z} J_{x}\comm{\Sz_{i}\Sz_{i+1}}{\Sx_{k}\Sx_{k+1}} + J_{z} J_{y}\comm{\Sz_{i}\Sz_{i+1}}{\Sy_{k}\Sy_{k+1}}  
\end{align*}
By inspection it becomes clear that out of six terms present, only three will need to be directly evaluated, as commutators of the form
\(\comm{A}{B}\) will differ from \(\comm{B}{A}\) by a sign and an index change.

\begin{align*}
    J_{x} J_{y}\comm{\Sx_{i}\Sx_{i+1}}{\Sy_{k}\Sy_{k+1}} =& J_x J_y \Big(\Sx_i\comm{\Sx_{i+1}}{\Sy_{k}\Sy_{k+1}} + \comm{\Sx_{i}}{\Sy_{k}\Sy_{k+1}} \Sx_{i+1} \Big)\\
    =& J_x J_y \Big( \Sx_i \left( \Sy_k \comm{\Sx_{i+1}}{\Sy_{k+1}} + \comm{\Sx_{i+1}}{\Sy_{k}} \Sy_{k+1}\right) + 
    \left( \Sy_k \comm{\Sx_{i}}{\Sy_{k+1}} + \comm{\Sx_{i}}{\Sy_{k}} \Sy_{k+1} \right)\Sx_{i+1} \Big)\\
    =& i J_x J_y \Big( \delta_{i+1,k+1} \Sx_i \Sy_k \Sz_{i+1} + \delta_{i+1,k} \Sx_i \Sz_{i+1} \Sy_{k+1} + \delta_{i,k+1} \Sy_k \Sz_i \Sx_{i+1} +\delta_{i,k} \Sz_i \Sy_{k+1} \Sx_{i+1}\Big)
\end{align*}
Carrying out the calculation of remaining two non-trivial commutators, we arrive at the following equations:
\begin{align*}
    J_{z} J_{x}\comm{\Sz_{i}\Sz_{i+1}}{\Sx_{k}\Sx_{k+1}} =& i J_z J_x \Big( \delta_{i+1,k+1} \Sx_k \Sz_i \Sy_{k+1} + \delta_{i+1,k} \Sz_i \Sy_{k} \Sx_{k+1} + \delta_{i,k+1} \Sx_k \Sy_{k+1} \Sz_{i+1} +\delta_{i,k} \Sy_k \Sz_{i+1} \Sx_{k+1}\Big)\\ 
    J_{y} J_{z}\comm{\Sy_{i}\Sy_{i+1}}{\Sz_{k}\Sz_{k+1}} =& i J_y J_z \Big( \delta_{i+1,k+1} \Sy_i \Sz_k \Sx_{i+1} + \delta_{i,k+1} \Sz_k \Sx_{i} \Sy_{i+1} + \delta_{i+1,k} \Sy_i \Sx_{i+1} \Sz_{k+1} +\delta_{i,k} \Sx_i \Sz_{k+1} \Sy_{i+1}\Big) 
\end{align*}
Next step requires us to evaluate the sum over lattice sites to get rid of the Kronecker \(\delta{}\)'s. As before, one of the three parts of calculations is provided in full detail:
\begin{align*}
&i\sum_{k=1}^{L}  J_{x} J_{y}\comm{\Sx_{i}\Sx_{i+1}}{\Sy_{k}\Sy_{k+1}} + i\sum_{k=1}^{L} J_{x} J_{y}\comm{\Sy_{i}\Sy_{i+1}}{\Sx_{k}\Sx_{k+1}} = \\
& - J_x J_y \Big( \cancel{\Sx_i \Sy_i \Sz_{i+1}} +  \Sx_i \Sz_{i+1} \Sy_{i+2} + \Sy_{i-1} \Sz_i \Sx_{i+1} +\bcancel{\Sz_i \Sy_{i+1} \Sx_{i+1}}\Big)\\
&+  J_x J_y \Big(  \cancel{\Sx_i \Sy_i \Sz_{i+1}} + \Sy_i \Sz_{i+1} \Sx_{i+2} + \Sx_{i-1} \Sz_i \Sy_{i+1} + \bcancel{\Sz_i \Sy_{i+1} \Sx_{i+1}}\Big)\\
&= J_x J_y \Big( \Sy_i \Sz_{i+1} \Sx_{i+2} - \Sx_i \Sz_{i+1} \Sy_{i+1} - \left(\Sy_{i-1}\Sz_i \Sx_{i+1} - \Sx_{i-1}\Sz_{i}\Sy_{i+1}\right) \Big)\\
\\
&i\sum_{k=1}^{L}  J_{x} J_{z}\comm{\Sx_{i}\Sx_{i+1}}{\Sz_{k}\Sz_{k+1}} + i\sum_{k=1}^{L} J_{x} J_{z}\comm{\Sz_{i}\Sz_{i+1}}{\Sx_{k}\Sx_{k+1}} = \\
&= J_x J_z \Big( \Sx_i \Sy_{i+1} \Sz_{i+2} - \Sz_i \Sy_{i+1} \Sx_{i+2} - \left(\Sx_{i-1}\Sy_i \Sz_{i+1} - \Sz_{i-1}\Sy_{i}\Sx_{i+1}\right) \Big)\\
\\
&i\sum_{k=1}^{L}  J_{y} J_{z}\comm{\Sy_{i}\Sy_{i+1}}{\Sz_{k}\Sz_{k+1}} + i\sum_{k=1}^{L} J_{y} J_{z}\comm{\Sz_{i}\Sz_{i+1}}{\Sy_{k}\Sy_{k+1}} = \\
&= J_y J_z \Big( \Sz_i \Sx_{i+1} \Sy_{i+2} - \Sy_i \Sx_{i+1} \Sz_{i+2} - \left(\Sz_{i-1}\Sx_i \Sy_{i+1} - \Sy_{i-1}\Sx_{i}\Sz_{i+1}\right) \Big)\\
\end{align*}
What now remains is to collect these parts and see that we finally arrive at the equation for the local energy current:
\begin{align}
    j_i^E &= J_x J_y \left(\Sy_{i-1}\Sz_i \Sx_{i+1} - \Sx_{i-1}\Sz_{i}\Sy_{i+1}\right) \nonumber\\
    &+ J_x J_z \left(\Sx_{i-1}\Sy_i \Sz_{i+1} - \Sz_{i-1}\Sy_{i}\Sx_{i+1}\right) \nonumber \\
    &+ J_y J_z \left(\Sz_{i-1}\Sx_i \Sy_{i+1} - \Sy_{i-1}\Sx_{i}\Sz_{i+1}\right) \nonumber \\
    &= J_x J_y \left(\Sy_{i-1}\Sz_i \Sx_{i+1} - \Sx_{i-1}\Sz_{i}\Sy_{i+1}\right) + \text{cyclic permutations of } \left(x,y,z\right)
    \label{eq:local energy current}
\end{align}
which is precisely the expression from~\textcite{Zotos1997}. Obtaining the energy current operator is now simply the matter of summing over all the lattice sites:
\begin{equation}
    J^E = \sum_{i=1}^L j_i^E
    \label{eq:energy current}
\end{equation}

However, in this work we are interested in the XXZ model with the Hamiltonian~\eqref{eq:HXXZ}. To this end,
we need to set \(J_x, J_z = 2J\), \(J_z = \Delta\) and substitute \(\Sx_i = \frac{\Sp_i + \Sm_i}{2}\), \(\Sy_i = \frac{\Sp_i-\Sm_i}{2i}\).
	\cleardoublepage{}

	\chapter{scratchpad}
\thispagestyle{chapterBeginStyle}

In this work we will describe \ldots~\autocite{Mierzejewski2015a}



\noindent Quantity that is plotted in Figures~4--14:
\begin{itemize}
    \item 
    With extrapolation to thermodynamic limit:
    \begin{equation}
        R_l(\tau,\alpha) = \frac{\lambda_l(L\rightarrow \infty,\tau,\alpha)}
        {\lambda_l(L\rightarrow \infty,\tau \rightarrow \infty,\alpha=0)}
        \label{eq:R1 extrap}
    \end{equation}
    \item
    Without extrapolation to thermodynamic limit:
    \begin{equation}
        R^L_l(\tau,\alpha) = \frac{\lambda_l(L,\tau,\alpha)}
        {\lambda_l(L,\tau \rightarrow \infty,\alpha=0)}
        \label{eq:R1 no extrap}
    \end{equation}
\end{itemize}


\noindent Energy current in integrable XXZ model: 
\begin{align}
    J^E = \sum_i^L i\left[\beta_1 \left(\Sm_{i-1}\Sz_{i}\Sp_{i+1}\right) + \beta_2 \left(\Sz_{i-1}\Sp_{i}\Sm_{i+1}+\Sp_{i-1}\Sm_{i}\Sz_{i+1}\right)\right] + \text{H.c.}
\end{align}


\noindent Spin LIOM for \(\Delta=1.0\):
\begin{equation}
    \hat{O}_{1} =  \beta_3\sum_{i=1}^{L} \Sp_{i} + \text{H.c.}
\end{equation}

\noindent Spin QLIOM for \(\Delta=-0.5\):
\begin{equation}
    \hat{O}_{1} =\beta_4 \sum_{i=1}^{L} \Bigl( \Sp_{i}\Sp_{i+1}\Sp_{i+2}\Bigr) + \text{H.c.}
\end{equation}

\begin{equation*}
    \comm{J^E}{\Sz_{tot}} = 0 \text{   where   } \Sz_{tot} = \sum_i \Sz_i
  \end{equation*}
  
  
\section{Operators with largest stiffness}

In this section we list operators with leading eigenvalues.\\

\noindent Fermions with \(m = 3\):
\begin{itemize}
    \item \( \displaystyle \begin{aligned}[t]
        \hat{O}_{max} = & \sum_{i=1}^{L} \Bigl(\alpha_1(\Sp_{i}\Sp_{i+1}\Id_{i+2})+\alpha_2(\Sp_{i}\Sz_{i+1}\Sp_{i+2})\Bigr) \text{ for even \(L\) and \(\Delta = \pm 1.0\)}
    \end{aligned} \)
    \item \(\displaystyle \begin{aligned}[t] 
        \hat{O}_{max} =& \sum_{i=1}^{L} \Bigl(\alpha_1(\Sp_{i}\Sp_{i+1}\Id_{i+2})+ 
        \alpha_2(\Sp_{i}\Id_{i+1}\Sp_{i+2})+ \alpha_3(\Sz_{i}\Sp_{i+1}\Sp_{i+2})+
        \alpha_4(\Sp_{i}\Sz_{i+1}\Sp_{i+2})+\\ 
        &\alpha_5(\Sp_{i}\Sp_{i+1}\Sz_{i+2})\Bigr) \text{ for odd \(L\) and \(\Delta = \pm 1.0\)}
    \end{aligned}\)
     
     
     \item \( \displaystyle \begin{aligned}[t]
         \hat{O}_{max} =& \sum_{i=1}^{L} \Bigl(\alpha_1(\Sp_{i}\Sp_{i+1}\Id_{i+2})+\alpha_2(\Sp_{i}\Sz_{i+1}\Sp_{i+2})\Bigr)
         \text{ for \(L = 8,12\) and \(\Delta = \pm 0.5\)}
     \end{aligned}\)
      
     \item \( \displaystyle \begin{aligned}[t]
         \hat{O}_{max} =& \sum_{i=1}^{L}{(-1)}^{i} \alpha_1\Bigl(\Sp_{i}\Sp_{i+1}\Id_{i+2}\Bigr) 
         \text{ for \(L = 10,14\) and \(\Delta = \pm 0.5\)}
     \end{aligned}\)

     \item \( \displaystyle \begin{aligned}[t]
         \hat{O}_{max} =& \sum_{i=1}^{L} \Bigl(\alpha_1(\Sp_{i}\Sp_{i+1}\Id_{i+2}) + \alpha_2(\Sp_{i}\Id_{i+1}\Sp_{i+1})+
         \alpha_3(\Sz_{i}\Sp_{i+1}\Sp_{i+2})+ \alpha_4(\Sp_{i}\Sz_{i+1}\Sp_{i+2})+\\
         &\alpha_5(\Sp_{i}\Sp_{i+1}\Sz_{i+2})\Bigr)\text{ for \(L = 9,11,13\) and \(\Delta = \pm 0.5\)}
     \end{aligned}\)
\end{itemize}
\noindent Fermions with \(m = 4\):
\begin{itemize}
    \item \(\hat{O}_{max}=\sum_{i=1}^{L}  \)   
\end{itemize}


\noindent Spins with \(m = 3\):
\begin{itemize}
    \item \( \displaystyle \begin{aligned}[t]
        \hat{O}_{max} =& \sum_{i=1}^{L} \alpha_1\Bigl( \Sp_{i}\Id_{i+1}\Id_{i+2}\Bigr)
        \text{ for all \(L\) and \(\Delta = 1.0\)}
    \end{aligned} \)
    \(\alpha_1 = \pm\frac{1}{\sqrt{L}}\)
    \item \( \displaystyle \begin{aligned}[t]
        \hat{O}_{max} =& \sum_{i=1}^{L}{(-1)}^{i} \alpha_1\Bigl( \Sp_{i}\Id_{i+1}\Id_{i+2}\Bigr)
        \text{ for even \(L\) and \(\Delta = -1.0\)}
    \end{aligned} \)
    \(\alpha_1 = \pm\frac{1}{\sqrt{L}}\)
    \item \( \displaystyle \begin{aligned}[t]
        \hat{O}_{max} =& \sum_{i=1}^{L} \Bigl( \alpha_1(\Sp_{i}\Id_{i+1}\Id_{i+2}) + \alpha2( \Sm_{i} \Sp_{i+1}\Sp_{i+2}) +
        \alpha_3 ( \Sz_{i}\Sz_{i+1}\Sp_{i+2}) + \alpha_4 (\Sp_{i}\Sm_{i+1}\Sp_{i+2}) + \\
        & \alpha_5 (\Sm_{i}\Sm_{i+1}\Sp_{i+2}) + \alpha_6 (\Sz_{i}\Sp_{i+1}\Sz_{i}) + \alpha_6 (\Sp_{i}\Sz_{i+1}\Sz_{i+2})
             \Bigr) \text{ for odd \(L\) and \(\Delta = -1.0\)}
    \end{aligned} \)
    \item \( \displaystyle \begin{aligned}[t]
        \hat{O}_{max} =& \sum_{i=1}^{L} \alpha_1\Bigl( \Sp_{i}\Sp_{i+1}\Sp_{i+2}\Bigr)
        \text{ for all \(L\) and \(\Delta = -0.5\)}
    \end{aligned} \)
    \item \( \displaystyle \begin{aligned}[t]
        \hat{O}_{max} =& \sum_{i=1}^{L} {(-1)}^{i}\alpha_1\Bigl( \Sp_{i}\Id_{i+1}\Id_{i+2}\Bigr)
        \text{ for even \(L\) and \(\Delta = 0.5\)}
    \end{aligned} \)
    \item \( \displaystyle \begin{aligned}[t]
        \hat{O}_{max} =& \sum_{i=1}^{L} {(-1)}^{i}\Bigl( \alpha_1(\Sp_{i}\Id_{i+1}\Id_{i+2}) + \alpha2( \Sm_{i} \Sp_{i+1}\Sp_{i+2}) +
        \alpha_3 ( \Sz_{i}\Sz_{i+1}\Sp_{i+2}) + \alpha_4 (\Sp_{i}\Sm_{i+1}\Sp_{i+2}) + \\
        & \alpha_5 (\Sm_{i}\Sm_{i+1}\Sp_{i+2}) + \alpha_6 (\Sz_{i}\Sp_{i+1}\Sz_{i}) + \alpha_6 (\Sp_{i}\Sz_{i+1}\Sz_{i+2})
             \Bigr) \text{ for odd \(L\) and \(\Delta = 0.5\)}
    \end{aligned} \)
    
\end{itemize}

Konwencja w kodzie: \(0 \longleftrightarrow \mathcal{I} \), \(1 \longleftrightarrow S^{+}\), \(2 \longleftrightarrow\ S^{z}\),
 \(3 \longleftrightarrow S^{-}\) 


\noindent Spin operators supported on up to \(m = 3\) sites:
\begin{itemize}
    \item \( \displaystyle \begin{aligned}[t]
        \hat{O}_{1} =&  \alpha_1\sum_{i=1}^{L} \Sp_{i} + \text{H.c.},
        \text{ for all \(L\) and \(\Delta = 1.0\)}
    \end{aligned} \)
    \(\alpha_1 = \pm\frac{1}{\sqrt{L}}\)
    \item \( \displaystyle \begin{aligned}[t]
        \hat{O}_{1} =& \alpha_1\sum_{i=1}^{L}{(-1)}^{i}\Sp_{i} + \text{H.c.},
        \text{ for even \(L\) and \(\Delta = -1.0\)}
    \end{aligned} \)
    \(\alpha_1 = \pm\frac{1}{\sqrt{L}}\)
    \item \( \displaystyle \begin{aligned}[t]
        \hat{O}_{1} =& \sum_{i=1}^{L} \Bigl( \alpha_1(\Sp_{i}) + \alpha_2( \Sm_{i} \Sp_{i+1}\Sp_{i+2}) +
        \alpha_3 ( \Sz_{i}\Sz_{i+1}\Sp_{i+2}) + \alpha_4 (\Sp_{i}\Sm_{i+1}\Sp_{i+2}) + \\
        & \alpha_5 (\Sm_{i}\Sm_{i+1}\Sp_{i+2}) + \alpha_6 (\Sz_{i}\Sp_{i+1}\Sz_{i}) + \alpha_7 (\Sp_{i}\Sz_{i+1}\Sz_{i+2})
             \Bigr) + \text{H.c.}, \text{ for odd \(L\) and \(\Delta = -1.0\)}
    \end{aligned} \)
    \item \( \displaystyle \begin{aligned}[t]
        \hat{O}_{1} =&\alpha_1 \sum_{i=1}^{L} \Bigl( \Sp_{i}\Sp_{i+1}\Sp_{i+2}\Bigr) + \text{H.c.,}
        \text{ for all \(L\) and \(\Delta = -0.5\)}
    \end{aligned} \)
    \(\alpha_1 = \pm\frac{1}{\sqrt{L}}\)
    \item \( \displaystyle \begin{aligned}[t]
        \hat{O}_{1} =& \alpha_1\sum_{i=1}^{L} {(-1)}^{i}\Bigl( \Sp_{i}\Sp_{i+1}\Sp_{i+2}\Bigr)
       + \text{H.c., for even \(L\) and \(\Delta = 0.5\)}
    \end{aligned} \)
    \item \( \displaystyle \begin{aligned}[t]
        \hat{O}_{1} =& \sum_{i=1}^{L} \Bigl( \alpha_1(\Sp_{i}) + \alpha_2( \Sm_{i} \Sp_{i+1}\Sp_{i+2}) +
        \alpha_3 ( \Sz_{i}\Sz_{i+1}\Sp_{i+2}) + \alpha_4 (\Sp_{i}\Sm_{i+1}\Sp_{i+2}) + \\
        & \alpha_5 (\Sm_{i}\Sm_{i+1}\Sp_{i+2}) + \alpha_6 (\Sz_{i}\Sp_{i+1}\Sz_{i}) + \alpha_7 (\Sp_{i}\Sz_{i+1}\Sz_{i+2})
             \Bigr) + \text{H.c.,} \text{ for odd \(L\) and \(\Delta = 0.5\)}
    \end{aligned} \)
\end{itemize}

\noindent Fermion operators supported on up to \(m = 3\) sites:
\begin{itemize}
    \item \( \displaystyle \begin{aligned}[t]
        \hat{O}_{1} = & \sum_{i=1}^{L} \Bigl(\alpha_1(\Sp_{i}\Sp_{i+1})+\alpha_2(\Sp_{i}\Sz_{i+1}\Sp_{i+2})\Bigr)+ \text{H.c., for even \(L\) and \(\Delta = \pm 1.0\)}
    \end{aligned} \)
    \item \(\displaystyle \begin{aligned}[t] 
        \hat{O}_{1} =& \sum_{i=1}^{L} \Bigl(\alpha_1(\Sp_{i}\Sp_{i+1})+ 
        \alpha_2(\Sp_{i}\Id_{i+1}\Sp_{i+2})+ \alpha_3(\Sz_{i}\Sp_{i+1}\Sp_{i+2})+
        \alpha_4(\Sp_{i}\Sz_{i+1}\Sp_{i+2})+\\ 
        &\alpha_5(\Sp_{i}\Sp_{i+1}\Sz_{i+2})\Bigr) +\text{H.c., for odd \(L\) and \(\Delta = \pm 1.0\)}
    \end{aligned}\)
     
     
     \item \( \displaystyle \begin{aligned}[t]
         \hat{O}_{1} =& \sum_{i=1}^{L} \Bigl(\alpha_1(\Sp_{i}\Sp_{i+1})+\alpha_2(\Sp_{i}\Sz_{i+1}\Sp_{i+2})\Bigr)
         + \text{H.c., for \(L = 8,12\) and \(\Delta = \pm 0.5\)}
     \end{aligned}\)
      
     \item \( \displaystyle \begin{aligned}[t]
         \hat{O}_{1} =&  \alpha_1 \sum_{i=1}^{L}{(-1)}^{i}\Bigl(\Sp_{i}\Sp_{i+1}\Bigr) 
         +\text{H.c., for \(L = 10,14\) and \(\Delta = \pm 0.5\)}
     \end{aligned}\)

     \item \( \displaystyle \begin{aligned}[t]
         \hat{O}_{1} =& \sum_{i=1}^{L} \Bigl(\alpha_1(\Sp_{i}\Sp_{i+1}) + \alpha_2(\Sp_{i}\Id_{i+1}\Sp_{i+1})+
         \alpha_3(\Sz_{i}\Sp_{i+1}\Sp_{i+2})+ \alpha_4(\Sp_{i}\Sz_{i+1}\Sp_{i+2})+\\
         &\alpha_5(\Sp_{i}\Sp_{i+1}\Sz_{i+2})\Bigr)+\text{H.c., for \(L = 9,11,13\) and \(\Delta = \pm 0.5\)}
     \end{aligned}\)
\end{itemize}
	\cleardoublepage{}
	
	%%%%%%%%%%%%%%%%%%%%%%%%%%%%%%%%%%%%%%%%%%%%%%%%%%%%%%%%%%%%%%%%%%%%%%%%%%%%%%
	%%%%%%%%%%%%%%%%%%%%%%%%%%%%%%% BIBLIOGRAFIA %%%%%%%%%%%%%%%%%%%%%%%%%%%%%%%%%
	%%%%%%%%%%%%%%%%%%%%%%%%%%%%%%%%%%%%%%%%%%%%%%%%%%%%%%%%%%%%%%%%%%%%%%%%%%%%%%

	\pagestyle{bibliographyStyle}
	\printbibliography{}
	\thispagestyle{chapterBeginStyle}
        \addcontentsline{toc}{chapter}{Bibliography}
	\cleardoublepage{}
	
	%%%%%%%%%%%%%%%%%%%%%%%%%%%%%%%%%%%%%%%%%%%%%%%%%%%%%%%%%%%%%%%%%%%%%%%%%%%%%%
	%%%%%%%%%%%%%%%%%%%%%%%%%%%%%%%%% DODATKI %%%%%%%%%%%%%%%%%%%%%%%%%%%%%%%%%%%%
	%%%%%%%%%%%%%%%%%%%%%%%%%%%%%%%%%%%%%%%%%%%%%%%%%%%%%%%%%%%%%%%%%%%%%%%%%%%%%%
	
	% \appendix
	% \pagestyle{appendixStyle}
	
	% \input{dodatek1}
	% \cleardoublepage

\end{document}

